\chapter{Planificación y metodología}\label{cap:planificacion}

Este capítulo describe cómo se ha organizado el trabajo para abordar un problema
de la magnitud del cloud lock-in en el marco de un Trabajo de Fin de Grado. Se
presenta la estrategia de desarrollo adoptada, la metodología y las herramientas
utilizadas, la planificación temporal y una estimación del presupuesto del proyecto.

\section{Estrategia de desarrollo}

El problema que aborda CloudPorter es amplio y sus dimensiones —abstracción de
infraestructura, estimación de costes, seguridad— no pueden desarrollarse de forma
simultánea en el marco de un Trabajo de Fin de Grado realizado por una sola persona.
La estrategia planteada consiste en construir de dentro hacia afuera: primero el
núcleo que hace viable el enfoque y, sobre él, las capas que lo enriquecen.

El núcleo es el flujo agnóstico de extremo a extremo de describir la infraestructura en
el Manifest, traducirla a las plantillas de un proveedor cloud y desplegarla. Sin esa
base funcionando, el resto del sistema carece de sentido, por lo que es lo primero
que se aborda. El soporte a un segundo proveedor se contempla como la prueba de que
el enfoque es realmente portable.

Sobre ese núcleo se prevén las capas que aportan valor sin necesidad de
infraestructura real, por operar sobre el propio Manifest: la estimación comparativa
de costes entre proveedores y la auditoría de seguridad, que detecta configuraciones
inseguras antes del despliegue. Junto a ellas, la cobertura de tipos de recurso
adicionales, como las bases de datos, amplía lo que el sistema puede representar.

La interfaz de línea de comandos es el punto de entrada que expone cada una de estas
funcionalidades a medida que se incorporan.

En cuanto a la elección de proveedores, se ha optado por \textbf{AWS} y
\textbf{Microsoft Azure} como plataformas objetivo de la versión inicial. AWS es
el líder indiscutible del mercado cloud y el punto de referencia para cualquier
herramienta de infraestructura; Azure es el segundo proveedor por cuota de mercado
y el más presente en entornos empresariales, precisamente el contexto donde el
lock-in tiene mayor impacto organizacional y económico \cite{flexera2026}. El
soporte a dos proveedores es suficiente para demostrar la viabilidad del enfoque
agnóstico.

\section{Metodología}

En esta sección se describen primero el marco de gestión del trabajo, Kanban, y el ciclo
que recorre cada unidad de trabajo; a continuación, las decisiones técnicas de fondo y las
herramientas que dan soporte al proceso: el motor de infraestructura como código, el lenguaje
de implementación, el entorno de desarrollo, los controles de calidad y la automatización.

\subsection{Kanban}

Se ha adoptado \textbf{Kanban} como metodología de desarrollo por su adecuación
al carácter exploratorio del proyecto. Al emerger el diseño durante el propio
desarrollo, no es posible anticipar cuánto trabajo conllevará cada ciclo ni cuánto
durará, por lo que comprometerse de antemano con un volumen fijo de tareas por
iteración —como exigen los sprints de Scrum— resultaría artificial. Kanban, en
cambio, propone un flujo continuo de trabajo sin ese compromiso previo, lo que permite
reordenar las prioridades en función de los resultados que se van obteniendo. A esta
razón se suma una de orden práctico: la dedicación al proyecto no pudo ser constante,
al compaginarse con un empleo y otras asignaturas, sino concentrada en ráfagas, un
patrón al que ese flujo continuo se adapta con naturalidad.

El tablero Kanban se gestiona íntegramente en \textbf{GitHub Projects}, integrado
con el repositorio del código. Cada unidad de trabajo se representa como un
\textit{issue} de GitHub, que recorre las columnas del tablero desde su definición
hasta su cierre. Las columnas del tablero son:

\begin{enumerate}
    \item \textbf{Backlog:} ideas y tareas identificadas, pendientes de priorizar.
    \item \textbf{Ready:} tarea priorizada y lista para ser iniciada.
    \item \textbf{In progress:} tarea en desarrollo activo.
    \item \textbf{In review:} pull request abierto, pendiente de revisión y checks de CI.
    \item \textbf{Done:} tarea completada y mergeada en \texttt{main}.
\end{enumerate}

Cada issue da lugar a una rama de trabajo con el esquema
\texttt{feat/<número>-<descripción>}, que se integra en \texttt{main} mediante
un pull request una vez superados los checks de CI. La rama principal está
protegida y no recibe cambios directos.

Los mensajes de commit siguen la especificación \textbf{Conventional
Commits}~\cite{conventionalCommits}, que establece un formato estructurado,
\texttt{tipo(ámbito): descripción}, legible tanto por personas como por
herramientas de automatización.

Los issues de mayor entidad, o grupos de issues relacionados, dan lugar a ciclos
de desarrollo documentados en el capítulo~\ref{cap:desarrollo}, siguiendo la
estructura de objetivo de la iteración, decisiones de diseño, implementación y
resultado.

\subsection{Ciclo de desarrollo}

El ciclo de desarrollo para cada iteración sigue el siguiente flujo:

\begin{enumerate}
    \item \textbf{Definición del issue:} se describe el objetivo, el contexto y
    los criterios de aceptación en el issue de GitHub.

    \item \textbf{Diseño:} antes de implementar, se toman y documentan las
    decisiones de diseño relevantes: qué se construye, cómo y por qué.

    \item \textbf{Implementación:} esta fase se hace con el apoyo
    de \textbf{Claude Code}\cite{claudeCode}, un agente de programación
    cuyo uso se documenta explícitamente como parte de la metodología por su
    relevancia en el contexto actual del desarrollo de software.

    \item \textbf{Calidad local:} antes de subir el código se aplican controles
    automáticos en el entorno local —formato, tipos y tests— que actúan como
    primera barrera de calidad. El objetivo es que ningún cambio salga del
    entorno de desarrollo sin cumplir los estándares definidos.

    \item \textbf{Pull request y revisión:} el trabajo se integra mediante un
    pull request sobre el que se ejecutan automáticamente los checks de
    integración continua como segunda barrera antes del merge. La revisión
    verifica además la coherencia del cambio con las decisiones de diseño
    de la iteración.

    \item \textbf{Documentación:} una vez cerrado el issue, la iteración queda
    recogida en el capítulo~\ref{cap:desarrollo} con su objetivo, las decisiones
    tomadas, la implementación realizada y el resultado obtenido.
\end{enumerate}

\subsection{Motor de infraestructura como código}\label{subsec:decisiones}

Una de las decisiones transversales del proyecto es el uso de \textbf{OpenTofu}
\cite{opentofuDocs} como motor de generación y despliegue de infraestructura, en
lugar de \textbf{Terraform}. Aunque Terraform es la herramienta IaC más extendida,
en agosto de 2023 su empresa matriz, HashiCorp, cambió su licencia de
\textit{Mozilla Public License 2.0} (MPL-2.0, open source) a
\textit{Business Source License 1.1} (\hyperlink{bsl}{BSL}-1.1) \cite{terraformDocs}. Esta licencia
restringe el uso comercial a entidades que compitan directamente con HashiCorp,
pero no define con precisión qué constituye competencia, lo que genera una zona
gris legal para cualquier proyecto que construya herramientas de infraestructura
sobre Terraform.

Como respuesta de la comunidad open source, se creó \textbf{OpenTofu} bajo el
paraguas de la \textbf{Linux Foundation}, partiendo de la última versión de
Terraform con licencia MPL-2.0. OpenTofu mantiene compatibilidad con el ecosistema
Terraform y elimina la ambigüedad legal. Para CloudPorter, que genera y distribuye
plantillas de infraestructura como parte de su funcionalidad, la elección de
OpenTofu es la única opción que garantiza libertad de uso sin restricciones.

\subsection{Lenguaje de implementación}

En el plano del lenguaje de implementación, el proyecto se desarrolla íntegramente
en \textbf{Python}. La elección responde a tres factores: un ecosistema de
librerías maduro para el procesamiento de YAML y JSON, formatos centrales en
la especificación de infraestructura, una amplia disponibilidad de herramientas
para el desarrollo de interfaces de línea de comandos, y la familiaridad previa
del autor con el lenguaje, que reduce la fricción en un proyecto de alcance
acotado. La gestión de dependencias y entornos virtuales se realiza con
\textbf{uv} \cite{uvDocs}, una herramienta de nueva generación que sustituye
a pip y poetry con una resolución de dependencias más rápida y resultados
deterministas.

\subsection{Entorno de desarrollo}

El entorno de desarrollo se define mediante un \hyperlink{devcontainer}{\textbf{devcontainer}} \cite{devcontainerDocs}.
Esta decisión responde a dos objetivos: mantener las dependencias explícitas y
versionadas en un único lugar, y eliminar la fricción de configuración para
cualquier colaborador.

\subsection{Herramientas de calidad}

Las herramientas de calidad no han sido elecciones independientes, sino un ecosistema
diseñado para automatizar los estándares de mantenibilidad a nivel de commit. 
Antes de cada commit, \textbf{pre-commit} \cite{precommitDocs} actúa como guardián 
ejecutando dos comprobaciones: el formateador y linter \textbf{Ruff} \cite{ruffDocs}
y la verificación estática de tipos con \textbf{mypy} \cite{mypyDocs}. Ambas herramientas
garantizan que ningún cambio se integra sin cumplir el estándar de formato y
tipos, proporcionando feedback inmediato sin salir del flujo de trabajo local.

Los tests automáticos se implementan con \textbf{pytest} \cite{pytestDocs} y se
ejecutan en dos momentos. Antes de cada push, la suite actúa como última
comprobación local que evita subir código roto. En integración continua se ejecuta
de nuevo, esta vez con informe de cobertura y un objetivo mínimo del 80\%.

\subsection{Automatización}

El proyecto adopta una práctica de \hyperlink{cicd}{CI/CD} mediante
\textbf{GitHub Actions}, que ejecuta los workflows necesarios en cada push,
pull request o publicación de tag. Se han configurado cuatro workflows:

\begin{itemize}
    \item \textbf{CI} (\texttt{ci.yml}): se ejecuta en cada push y pull request
    sobre \texttt{main}. Instala dependencias con uv, ejecuta Ruff, mypy y
    pytest con informe de cobertura. En pull requests, publica automáticamente
    un comentario con el resumen de cobertura mediante
    \texttt{python-coverage-comment-action}.

    \item \textbf{PR Conventions} (\texttt{pr-conventions.yml}): verifica que
    el título de cada pull request sigue la especificación Conventional Commits,
    rechazando la integración si no cumple el formato.

    \item \textbf{Sync PR from issue} (\texttt{sync-pr-from-issue.yml}): al
    abrir un pull request, extrae el número de issue del nombre de la rama,
    copia sus etiquetas al PR y añade automáticamente la referencia
    \texttt{Closes \#N} al cuerpo del PR para garantizar el cierre automático
    del issue al hacer merge.

    \item \textbf{Release} (\texttt{release.yml}): se activa al publicar un
    tag con formato \texttt{v*.*.*} sobre \texttt{main}. Construye el paquete
    con \texttt{uv build} y crea una GitHub Release con los artefactos
    adjuntos y las notas de cambios generadas automáticamente.
\end{itemize}

\subsection{Gestión de la memoria}

La documentación de este TFG se ha redactado con el apoyo de \textbf{Claude Code}
\cite{claudeCode}. Se configuraron \hyperlink{skill}{\textit{skills}} propias, que encapsulan el estilo 
de escritura académica y las convenciones LaTeX del documento, y un fichero \texttt{CLAUDE.md} 
que define el contexto del proyecto y el enfoque narrativo.

\section{Cronograma}

El proyecto se divide en dos fases principales, una de investigación y otra de
desarrollo, cuyo reparto temporal se resume en la figura~\ref{fig:gantt}.

La \textbf{fase de investigación} (febrero–abril de 2026) no siguió un recorrido
lineal, sino un proceso exploratorio organizado en \textit{spikes}: investigaciones
acotadas en el tiempo, cada una orientada a reducir una incertidumbre concreta. Se
sucedieron, con solapes entre ellas, las siguientes:

\begin{itemize}
    \item \textbf{S1. Detección del problema} (2--15 feb): las fricciones que el autor
    encuentra en su propia experiencia con la nube, que motivan preguntarse si son un
    caso aislado o un problema generalizado.

    \item \textbf{S2. Investigación del problema} (9 feb -- 15 mar): estudio de hasta
    qué punto el cloud lock-in afecta a otros y de sus causas y dimensiones.

    \item \textbf{S3. Exploración de soluciones} (23 feb -- 22 mar): revisión, en
    paralelo con S2, de las herramientas existentes, que confirma que ninguna resolvía
    el problema de forma satisfactoria.

    \item \textbf{S4. Conceptualización y alcance} (16 mar -- 5 abr): definición de la
    propuesta y delimitación de lo que entra en el trabajo.

    \item \textbf{S5. Metodología y entorno} (30 mar -- 19 abr): elección de la
    metodología de trabajo y montaje del entorno de desarrollo.
\end{itemize}

La formación en las tecnologías necesarias (2 feb -- 22 mar) acompañó de forma
transversal a todo el proceso, en parte previa al inicio formal del trabajo. Esta fase
culmina con la redacción de los capítulos 1 a 5.

La \textbf{fase de desarrollo} (abril–junio de 2026) implementa CloudPorter de forma
iterativa a lo largo de las ocho iteraciones documentadas en el
capítulo~\ref{cap:desarrollo}, con redacción simultánea de la memoria y una fase final
de evaluación y validación.

\begin{figure}[H]
\centering
\resizebox{\textwidth}{!}{%
\begin{ganttchart}[
    hgrid, vgrid,
    x unit=0.55cm,
    y unit chart=0.6cm,
    y unit title=0.7cm,
    title height=1,
    bar height=0.6,
    bar/.append style={fill=blue!45, draw=blue!60},
    group/.append style={fill=black!75, draw=black!75},
    group right shift=0,
    group top shift=0.1,
    group height=0.5,
    bar label font=\small,
    group label font=\small\bfseries,
    title label font=\small,
    milestone/.append style={fill=red!70}
]{1}{20}
  \gantttitle{2026}{20} \\
  \gantttitle{Febrero}{4}
  \gantttitle{Marzo}{4}
  \gantttitle{Abril}{5}
  \gantttitle{Mayo}{4}
  \gantttitle{Junio}{3} \\
  \ganttgroup{Investigación}{1}{11} \\
  \ganttbar{Formación (AWS, OpenTofu)}{1}{7} \\
  \ganttbar{S1 · Detección del problema}{1}{2} \\
  \ganttbar{S2 · Investigación del problema}{2}{6} \\
  \ganttbar{S3 · Exploración de soluciones}{4}{7} \\
  \ganttbar{S4 · Conceptualización y alcance}{7}{9} \\
  \ganttbar{S5 · Metodología y entorno}{9}{11} \\
  \ganttgroup{Desarrollo}{13}{20} \\
  \ganttbar{Iteración 1 · Infra + manifiesto}{13}{14} \\
  \ganttbar{Iteración 2 · Traducción AWS}{14}{14} \\
  \ganttbar{Iteración 3 · Comando deploy}{15}{15} \\
  \ganttbar{Iteración 4 · Soporte Azure}{16}{16} \\
  \ganttbar{Iteración 5 · Estimación de costes}{17}{17} \\
  \ganttbar{Iteración 6 · Base de datos}{17}{17} \\
  \ganttbar{Iteración 7 · Auditoría}{17}{17} \\
  \ganttbar{Iteración 8 · Demo de tres capas}{17}{17} \\
  \ganttbar{Redacción de la memoria}{5}{20} \\
  \ganttbar{Evaluación y validación}{18}{20}
\end{ganttchart}%
}
\caption{Cronograma del proyecto: diagrama de Gantt de las fases de investigación y desarrollo.}
\label{fig:gantt}
\end{figure}

\section{Presupuesto}

A continuación se presenta una estimación del coste total del proyecto, calculado como
si se tratase de un encargo profesional. La tabla~\ref{tab:presupuesto} desglosa los
conceptos considerados y su coste.

El coste de personal corresponde a una dedicación estimada de \textbf{326 horas} a una
tarifa de 16~€/hora, propia de un perfil de ingeniero junior con especialización en
entornos cloud. El cómputo es una estimación orientativa: dado el carácter exploratorio
del proyecto y su gestión mediante Kanban, el trabajo se desarrolló de forma discontinua,
en periodos de dedicación concentrada según la disponibilidad. Para aproximarla, se
registró el tiempo al inicio y al cierre de cada unidad de trabajo, cada issue de
desarrollo y cada spike de investigación, lo que arroja unas 326 horas: en torno a 190
de investigación, formación y documentación (capítulos 1 a 5) y 136 de desarrollo,
implementación y evaluación.

El resto de conceptos se estima de la forma siguiente:

\begin{itemize}
    \item \textbf{Equipamiento informático:} un ordenador de gama media, suficiente porque
    el despliegue se realiza en la nube y no exige hardware potente en local. Se estima un precio estimado de 1.200€.

    \item \textbf{Infraestructura cloud:} sin coste. Las pruebas en Azure y en AWS se
    cubrieron con los niveles gratuitos que ofrecen ambos proveedores, sin llegar a generar
    gasto.

    \item \textbf{Software:} nulo. Todas las herramientas empleadas son de código abierto o
    de uso gratuito —Python, OpenTofu, GitHub en su plan gratuito, uv, Ruff y pytest—, junto
    con Claude Code en su modalidad de uso académico.

    \item \textbf{Suministros:} conexión a internet y electricidad del equipo durante los
    meses de trabajo, a precios domésticos habituales.

    \item \textbf{Formación:} sin coste. Durante la investigación se emplearon cursos y
    recursos formativos gratuitos —tutoriales, documentación oficial y cursos en vídeo— y se
    obtuvo la certificación OpenTofu Foundation, de la Linux Foundation, también gratuita.
\end{itemize}

\begin{table}[H]
\centering
\begin{tabular}{lr}
\hline
\textbf{Concepto} & \textbf{Coste} \\
\hline
Personal & 5.216,00 € \\
Equipamiento informático & 1.200€ € \\
Software & 0,00 € \\
Infraestructura cloud & 0,00 € \\
Formación & 0,00 € \\
Suministros (internet y electricidad) & 300,00 € \\
\hline
\textbf{Total} & \textbf{6.716,00 €} \\
\hline
\end{tabular}
\caption{Desglose del presupuesto del proyecto.}
\label{tab:presupuesto}
\end{table}
